\section{Conclusions}\label{sec:conclusions}
%======================================================================

We presented a two-level adaptive LOD framework that certifies the error to
the exact Helmholtz solution.  The fine finite element space is used only
for local correctors and post-processing; it is not a reference problem.
The exact upper bound combines a compatibility-corrected equilibrated
majorant of the full residual with computable continuous-kernel defects.
The orthogonality of the continuous ideal trial and test spaces reduces the
stability perturbation to a product of the primal and adjoint defects, and
the required ideal inf--sup constant is bounded from the localized coarse
matrix.  Under the checked acceptance condition, the reliability constant
is independent of the wavenumber and all discretization parameters.

The adaptive realization separates inexpensive coarse marking from costly
exact certification.  Cache batches retain local factorizations while the
fine mesh and oversampling depth are fixed, but they do not alter the
mathematical target.  A failed certificate triggers either additional
oversampling or fine-mesh refinement, according to a discrete localization
diagnostic.  The construction is fully parallel at patch level and can
reuse the same mixed factors for the full residual and both total-defect
operators.

Several questions remain open.  In particular, local efficiency of the
fine-defect majorants, estimator reduction under fine refinement, and a
contraction or optimal-complexity theorem for the coupled adaptive loop are
not established here.  Sharper quotient minimization, certified iterative
eigenvalue solvers, and extensions to heterogeneous coefficients, higher
order elements, and scattering truncations are natural directions for
future work.  The multiresolution and super-localized Helmholtz variants in
\cite{HauckPeterseim2022,FreeseHauckPeterseim2024} provide further
opportunities to reduce the cost of the defect certification.
